% ============================================================================
% AVALIAÇÃO GERAL INTEGRADA (UNIDADES I A IV) — CAPÍTULOS 1 A 15
% ============================================================================

\cleardoublepage
\phantomsection
\addcontentsline{toc}{chapter}{Avaliação Geral Integrada (Unidades I a IV)}
\chapter*{Avaliação Geral Integrada (Unidades I a IV)}
\markboth{AVALIAÇÃO GERAL INTEGRADA}{AVALIAÇÃO GERAL INTEGRADA}

Esta avaliação final reúne \textbf{12 questões objetivas (A--E)} e \textbf{4 questões dissertativas integradoras}, cobrindo integralmente os Capítulos 1 a 15 das Unidades I, II, III e IV. Recomenda-se resolvê-la em condições simuladas de prova (duração sugerida: 2 horas, sem consulta ao gabarito) e registrar suas respostas na folha imprimível abaixo antes de conferir o espelho comentado.

\section*{Parte I: Questões Objetivas Integradas (12 Questões)}

\subsection*{Unidade I: Território, Estruturas de Mercado e Ecossistemas (Caps. 1 a 6)}

\begin{enumerate}[leftmargin=1.5em, itemsep=5pt]
  \item \textbf{(Caps. 1 e 3, Schumpeter, Milton Santos e Ecossistemas)} A crítica de Milton Santos ao uso técnico-informacional do território, articulada à visão evolucionária de Dosi e à abordagem de ecossistemas empreendedores de Stam, permite afirmar que uma política pública de inovação regional tende ao fracasso quando:
  \begin{alternativas}
    \item Fortalece a proximidade cognitiva moderada entre universidades locais e pequenas empresas do APL.
    \item Copia modelos institucionais de polos globais (isomorfismo) sem enraizamento na base produtiva e na governança local.
    \item Combina financiamento público paciente (Mazzucato) com capacitação de fornecedores regionais.
    \item Utiliza a Hélice Quíntupla para incorporar demandas socioambientais na agenda de P\&D.
    \item Articula Tecnologia Social e Indicações Geográficas para ampliar a apropriação local de renda.
  \end{alternativas}

  \item \textbf{(Cap. 2, Concentração de Mercado e HHI)} Em um município polo regional, quatro empresas concentram $100\%$ da oferta de serviços de logística refrigerada, com participações de mercado iguais a $50\%$, $25\%$, $15\%$ e $10\%$. O Índice Herfindahl-Hirschman ($\text{HHI}$) desse mercado e sua classificação usual são:
  \begin{alternativas}
    \item $\text{HHI} = 2.500$; mercado desconcentrado.
    \item $\text{HHI} = 3.125$; mercado moderadamente concentrado.
    \item $\text{HHI} = 3.450$; mercado altamente concentrado.
    \item $\text{HHI} = 4.000$; monopólio puro.
    \item $\text{HHI} = 1.850$; mercado perfeitamente competitivo.
  \end{alternativas}

  \item \textbf{(Caps. 4, 5 e 6, Hélice, Negócios de Impacto e Tecnologia Social)} Sobre a distinção entre empreendedorismo tradicional, negócios de impacto territorial e Tecnologia Social (Dagnino), assinale a afirmação correta:
  \begin{alternativas}
    \item A Tecnologia Social caracteriza-se por patentes fechadas de alto custo e gestão estritamente centralizada fora da comunidade.
    \item O fenômeno do \textit{mission drift} ocorre quando um negócio de impacto abandona sua finalidade socioambiental primária para atender apenas à pressão por retorno financeiro de curto prazo.
    \item A Hélice Tríplice exclui o Estado das políticas de CT\&I, restringindo-se a empresas, bancos e investidores anjo.
    \item Na Teoria da Mudança, ``insumos'' e ``impacto estrutural de longo prazo'' são sinônimos operacionais.
    \item Uma Indicação Geográfica (IG) impede a participação de cooperativas e associações de produtores familiares.
  \end{alternativas}
\end{enumerate}

\subsection*{Unidade II: Processo Empreendedor, Mercado e Fomento (Caps. 7 a 11)}

\begin{enumerate}[leftmargin=1.5em, itemsep=5pt, resume]
  \item \textbf{(Caps. 7 e 9, Effectuation, Lean Startup e Canvas)} De acordo com Sarasvathy (2001) e Osterwalder \& Pigneur (2010), a lógica \textit{effectuation} aplicada a estágios iniciais de alta incerteza tecnológica distingue-se da lógica \textit{causation} porque:
  \begin{alternativas}
    \item Parte de uma meta fixa e imutável de receita no Ano 5 antes de conhecer os recursos disponíveis.
    \item Elimina a necessidade de validar hipóteses com clientes reais (Blank/Ries), priorizando planos extensos em gabinete.
    \item Parte dos meios disponíveis (quem sou, o que sei, quem conheço), orienta-se pela perda aceitável e constrói parcerias estratégicas para cocriar o mercado.
    \item Proíbe a adaptação da proposta de valor ou o pivô no \textit{Business Model Canvas}.
    \item Aplica-se exclusivamente a grandes corporações estatais com monopólio legal.
  \end{alternativas}

  \item \textbf{(Cap. 8, Funil de Mercado TAM/SAM/SOM)} Uma \textit{agtech} baiana estima que existem $25.000$ estabelecimentos agropecuários potenciais no estado ($\text{TAM}$ em unidades). Desse total, $40\%$ possuem conectividade e perfil produtivo elegível à solução ($\text{SAM}$). A empresa projeta capturar $5\%$ desse mercado acessível em três anos ($\text{SOM}$), a um ticket médio anual de $\reais{2.400}$ por estabelecimento. A receita anual correspondente ao $\text{SOM}$ é de:
  \begin{alternativas}
    \item $\reais{600.000}$.
    \item $\reais{1.200.000}$.
    \item $\reais{2.400.000}$.
    \item $\reais{3.000.000}$.
    \item $\reais{12.000.000}$.
  \end{alternativas}

  \item \textbf{(Caps. 9, 10 e 11, Métricas de Tração, Fomento e Spin-offs)} Uma \textit{spin-off} acadêmica B2B cobra mensalidade de $\reais{250}$ por cliente, opera com margem de contribuição de $80\%$, apresenta taxa de cancelamento mensal (\textit{churn}) de $5\%$ e custo de aquisição de cliente ($\text{CAC}$) de $\reais{1.000}$. Considerando vida média do cliente igual a $1/\text{churn}$, o $\text{LTV}$ e a razão $\text{LTV/CAC}$ são, respectivamente:
  \begin{alternativas}
    \item $\text{LTV} = \reais{2.500}$ e $\text{LTV/CAC} = 2{,}5$.
    \item $\text{LTV} = \reais{3.200}$ e $\text{LTV/CAC} = 3{,}2$.
    \item $\text{LTV} = \reais{4.000}$ e $\text{LTV/CAC} = 4{,}0$.
    \item $\text{LTV} = \reais{5.000}$ e $\text{LTV/CAC} = 5{,}0$.
    \item $\text{LTV} = \reais{1.000}$ e $\text{LTV/CAC} = 1{,}0$.
  \end{alternativas}
\end{enumerate}

\subsection*{Unidade III: Projetos de Investimento, Localização e Valuation (Caps. 12 a 14)}

\begin{enumerate}[leftmargin=1.5em, itemsep=5pt, resume]
  \item \textbf{(Cap. 12, Matriz de Decisão Locacional)} Uma indústria de bioinsumos avalia três municípios (X, Y e Z) por meio de matriz ponderada com três critérios: C1, Proximidade de matéria-prima (peso $0{,}40$); C2, Qualificação técnica e ICT (peso $0{,}35$); C3, Incentivos e infraestrutura distrital (peso $0{,}25$). As notas atribuídas são: Município X $(8; 7; 9)$, Município Y $(9; 8; 5)$ e Município Z $(7; 9; 8)$. As pontuações finais e a alternativa vencedora são:
  \begin{alternativas}
    \item $S_X = 7{,}90$; $S_Y = 7{,}65$; $S_Z = 7{,}95$ $\rightarrow$ vence o Município Z.
    \item $S_X = 8{,}00$; $S_Y = 7{,}80$; $S_Z = 7{,}70$ $\rightarrow$ vence o Município X.
    \item $S_X = 7{,}65$; $S_Y = 8{,}10$; $S_Z = 7{,}95$ $\rightarrow$ vence o Município Y.
    \item $S_X = 7{,}90$; $S_Y = 7{,}65$; $S_Z = 7{,}85$ $\rightarrow$ vence o Município X.
    \item $S_X = 8{,}15$; $S_Y = 7{,}35$; $S_Z = 7{,}95$ $\rightarrow$ vence o Município X.
  \end{alternativas}

  \item \textbf{(Cap. 13, VPL e Payback Descontado)} Um projeto de modernização tecnológica exige investimento inicial de $\reais{100.000}$ e gera fluxos de caixa livres de $\reais{44.000}$ (Ano 1), $\reais{60.500}$ (Ano 2) e $\reais{66.550}$ (Ano 3), sob $\text{TMA}$ de $10\%$ ao ano. O $\text{VPL}$ e o \textit{Payback} Descontado desse projeto são:
  \begin{alternativas}
    \item $\text{VPL} = \reais{71.050{,}00}$ e $PB_d = 1{,}93\text{ ano}$.
    \item $\text{VPL} = \reais{40.000{,}00}$ e $PB_d = 2{,}20\text{ anos}$.
    \item $\text{VPL} = \reais{35.000{,}00}$ e $PB_d = 2{,}50\text{ anos}$.
    \item $\text{VPL} = \reais{40.000{,}00}$ e $PB_d = 2{,}80\text{ anos}$.
    \item $\text{VPL} = -\reais{10.000{,}00}$ e não há recuperação do capital.
  \end{alternativas}

  \item \textbf{(Cap. 14, WACC, Enterprise Value e Equity Value)} Uma empresa de base tecnológica possui estrutura de capital composta por $70\%$ de capital próprio ($K_e = 15\%$ a.a.) e $30\%$ de dívida ($K_d = 10\%$ a.a.), com alíquota de imposto de renda de $34\%$. Se o seu \textit{Enterprise Value} ($\text{EV}$) apurado por Fluxo de Caixa Descontado for de $\reais{1.500.000}$, com dívida bruta de $\reais{450.000}$ e caixa disponível de $\reais{120.000}$, então o $\text{WACC}$ e o \textit{Equity Value} valem, respectivamente:
  \begin{alternativas}
    \item $\text{WACC} = 13{,}50\%$ e $\text{Equity Value} = \reais{1.050.000}$.
    \item $\text{WACC} = 12{,}48\%$ e $\text{Equity Value} = \reais{1.170.000}$.
    \item $\text{WACC} = 12{,}48\%$ e $\text{Equity Value} = \reais{1.830.000}$.
    \item $\text{WACC} = 11{,}82\%$ e $\text{Equity Value} = \reais{1.170.000}$.
    \item $\text{WACC} = 12{,}00\%$ e $\text{Equity Value} = \reais{930.000}$.
  \end{alternativas}
\end{enumerate}

\subsection*{Unidade IV: Síntese Territorial e Seminário Integrador (Cap. 15)}

\begin{enumerate}[leftmargin=1.5em, itemsep=5pt, resume]
  \item \textbf{(Cap. 15, Articulação Epistemológica e Retenção Territorial)} Na perspectiva integradora do Seminário Final (AV3), um empreendimento tecnológico só promove desenvolvimento territorial efetivo quando:
  \begin{alternativas}
    \item Maximiza o $\text{VPL}$ privado por meio da importação integral de insumos e remessa total de dividendos para fora da região.
    \item Combina viabilidade econômico-financeira ($\text{VPL} > 0$, $\text{TIR} > \text{WACC}$) com elevado Índice de Retenção Territorial de Valor ($\text{IRV}$), encadeamento em APLs/SRIs e governança na Hélice Quádrupla/Quíntupla.
    \item Substitui a análise de fluxo de caixa por declarações qualitativas de boas intenções socioambientais.
    \item Ignora os limites do $\text{SOM}$ regional para adotar projeções globais padronizadas sem validação local.
    \item Evita qualquer interação com editais públicos regionais (FAPESB, BNB, FINEP) para não prestar contas territoriais.
  \end{alternativas}

  \item \textbf{(Cap. 15, Rubrica Oficial AV3 e Média Final)} Um estudante obteve $\text{AV1} = 7{,}5$ e $\text{AV2} = 8{,}0$. Na apresentação da AV3, recebeu $L_T = 9{,}0$, $E_M = 8{,}0$, $V_F = 8{,}4$ e $C_D = 8{,}0$. Aplicando os pesos oficiais da rubrica AV3 ($30\%$, $30\%$, $25\%$, $15\%$) e os pesos da disciplina ($30\%$ AV1, $20\%$ AV2, $50\%$ AV3), a nota $N_{\text{AV3}}$ e a Média Final ($M_F$) são:
  \begin{alternativas}
    \item $N_{\text{AV3}} = 8{,}35$ e $M_F = 8{,}02$.
    \item $N_{\text{AV3}} = 8{,}40$ e $M_F = 8{,}05$.
    \item $N_{\text{AV3}} = 8{,}40$ e $M_F = 7{,}95$.
    \item $N_{\text{AV3}} = 8{,}50$ e $M_F = 8{,}10$.
    \item $N_{\text{AV3}} = 8{,}20$ e $M_F = 7{,}95$.
  \end{alternativas}

  \item \textbf{(Cap. 15, Defesa Oral e Publicação Científica)} Em relação à defesa oral da AV3 e ao encaminhamento de resultados para publicação em fóruns como ENANPUR, ENANPAD e ENGEMA, assinale a alternativa correta:
  \begin{alternativas}
    \item A defesa oral dura $45$ minutos contínuos sem arguição, dispensando o limite de $10$ slides.
    \item O roteiro de $10$ slides em $20$ minutos deve equilibrar diagnóstico territorial, modelo de negócio, decisão locacional, viabilidade financeira e \textit{Valuation}, reservando $10$ minutos para arguição.
    \item Para publicar no ENANPUR ou ENANPAD, basta copiar o \textit{pitch} comercial sem pergunta de pesquisa nem referencial teórico.
    \item O ENGEMA aceita exclusivamente estudos sobre bolsas de valores internacionais, rejeitando bioeconomia e Tecnologia Social.
    \item A análise de sensibilidade é irrelevante em artigos científicos de planejamento territorial e projetos.
  \end{alternativas}
\end{enumerate}

\begin{tcolorbox}[colback=bgSoft, colframe=etIndigo, arc=3pt, boxrule=0.8pt, title={\bfseries Folha de Respostas Imprimível, Avaliação Geral Integrada (Parte I)}, coltitle=white, colbacktitle=etIndigo]
\textbf{Nome:} \rule{7.2cm}{0.4pt} \hfill \textbf{Matrícula:} \rule{3.0cm}{0.4pt} \hfill \textbf{Data:} \rule{2.2cm}{0.4pt}\\[6pt]
\renewcommand{\arraystretch}{1.25}
\centering
\footnotesize
\begin{tabularx}{\textwidth}{@{}c c c c c c | c c c c c c@{}}
  \toprule
  \rowcolor{bgGray}
  \textbf{Q.} & \textbf{A} & \textbf{B} & \textbf{C} & \textbf{D} & \textbf{E} & \textbf{Q.} & \textbf{A} & \textbf{B} & \textbf{C} & \textbf{D} & \textbf{E} \\
  \midrule
  \textbf{01} & $\bigcirc$ & $\bigcirc$ & $\bigcirc$ & $\bigcirc$ & $\bigcirc$ & \textbf{07} & $\bigcirc$ & $\bigcirc$ & $\bigcirc$ & $\bigcirc$ & $\bigcirc$ \\
  \textbf{02} & $\bigcirc$ & $\bigcirc$ & $\bigcirc$ & $\bigcirc$ & $\bigcirc$ & \textbf{08} & $\bigcirc$ & $\bigcirc$ & $\bigcirc$ & $\bigcirc$ & $\bigcirc$ \\
  \textbf{03} & $\bigcirc$ & $\bigcirc$ & $\bigcirc$ & $\bigcirc$ & $\bigcirc$ & \textbf{09} & $\bigcirc$ & $\bigcirc$ & $\bigcirc$ & $\bigcirc$ & $\bigcirc$ \\
  \textbf{04} & $\bigcirc$ & $\bigcirc$ & $\bigcirc$ & $\bigcirc$ & $\bigcirc$ & \textbf{10} & $\bigcirc$ & $\bigcirc$ & $\bigcirc$ & $\bigcirc$ & $\bigcirc$ \\
  \textbf{05} & $\bigcirc$ & $\bigcirc$ & $\bigcirc$ & $\bigcirc$ & $\bigcirc$ & \textbf{11} & $\bigcirc$ & $\bigcirc$ & $\bigcirc$ & $\bigcirc$ & $\bigcirc$ \\
  \textbf{06} & $\bigcirc$ & $\bigcirc$ & $\bigcirc$ & $\bigcirc$ & $\bigcirc$ & \textbf{12} & $\bigcirc$ & $\bigcirc$ & $\bigcirc$ & $\bigcirc$ & $\bigcirc$ \\
  \bottomrule
\end{tabularx}
\end{tcolorbox}

\section*{Parte II: Questões Dissertativas Integradoras (4 Questões)}

\begin{enumerate}[leftmargin=1.5em, itemsep=6pt]
  \item \textbf{(Unidade I, Território, Eficiência Coletiva e Hélice Quíntupla)} Explique por que a mera instalação física de um parque tecnológico ou incubadora em um município periférico pode reproduzir o isomorfismo institucional e gerar baixo impacto regional. Na sua resposta, mobilize os conceitos de eficiência coletiva (Schmitz), proximidades (Boschma) e Hélice Quádrupla/Quíntupla.
  \espacoResolucao[3.2cm]

  \item \textbf{(Unidade II, Oportunidade, PI e Escada do Capital)} Um grupo de pesquisa desenvolveu um bioprocesso para reaproveitamento de resíduos agroindustriais. Descreva como articular a lógica \textit{effectuation}, o funil $\text{TAM/SAM/SOM}$, a proteção de Propriedade Intelectual (Lei 9.279/1996 e Lei 13.243/2016) e a escada de fomento regional (FAPESB, FINEP, BNB, BNDES) para atravessar o ``Vale da Morte''.
  \espacoResolucao[3.2cm]

  \item \textbf{(Unidade III, Localização, Viabilidade e Valuation)} Justifique por que um projeto com $\text{TIR}$ superior à $\text{TMA}$ pode, ainda assim, destruir valor territorial ou enfrentar restrições de financiamento. Explique o papel complementar da Matriz Locacional Ponderada, do $\text{VPL}$, do \textit{Payback} Descontado e da diferença entre $\text{EV}$ e \textit{Equity Value}.
  \espacoResolucao[3.2cm]

  \item \textbf{(Unidade IV, Síntese AV3 e Agenda de Pesquisa)} Suponha que você esteja defendendo seu Projeto Final Integrador (AV3) em $20$ minutos e queira, em seguida, submetê-lo ao ENANPUR, ENANPAD ou ENGEMA. Indique como estruturar os argumentos nos $10$ slides para atender aos quatro critérios da Rubrica AV3 e quais adaptações metodológicas são indispensáveis para transformar o projeto em artigo científico.
  \espacoResolucao[3.2cm]
\end{enumerate}

\clearpage
\section*{Gabarito Integralmente Comentado e Diagnóstico Final}

\begin{boxGabarito}[Gabarito Comentado: Parte I (12 Questões Objetivas)]
\begin{enumerate}[leftmargin=1.5em, itemsep=4pt]
  \item \textbf{(B)} O isomorfismo institucional ocorre ao importar modelos exógenos desvinculados da estrutura produtiva, das proximidades locais e da governança territorial.
  \item \textbf{(C)} $\text{HHI} = 50^2 + 25^2 + 15^2 + 10^2 = 2.500 + 625 + 225 + 100 = \mathbf{3.450}$ ($> 2.500$: mercado altamente concentrado).
  \item \textbf{(B)} \textit{Mission drift} designa o desvio da missão socioambiental original sob pressão financeira ou de investidores.
  \item \textbf{(C)} Na lógica \textit{effectuation} (Sarasvathy), parte-se dos meios disponíveis, com foco na perda aceitável e na construção de alianças estratégicas.
  \item \textbf{(B)} $\text{SAM}_{\text{unidades}} = 25.000 \times 0{,}40 = 10.000$ estabelecimentos; $\text{SOM}_{\text{unidades}} = 10.000 \times 0{,}05 = 500$ estabelecimentos; Receita $\text{SOM} = 500 \times \reais{2.400} = \mathbf{\reais{1.200.000}}$.
  \item \textbf{(C)} Vida média $= \frac{1}{0{,}05} = 20$ meses. $\text{LTV} = 250 \times 0{,}80 \times 20 = \mathbf{\reais{4.000}}$; $\text{LTV/CAC} = \frac{4.000}{1.000} = \mathbf{4{,}0}$.
  \item \textbf{(A)} $S_X = 0{,}40(8)+0{,}35(7)+0{,}25(9) = 3{,}20+2{,}45+2{,}25 = 7{,}90$; $S_Y = 0{,}40(9)+0{,}35(8)+0{,}25(5) = 3{,}60+2{,}80+1{,}25 = 7{,}65$; $S_Z = 0{,}40(7)+0{,}35(9)+0{,}25(8) = 2{,}80+3{,}15+2{,}00 = \mathbf{7{,}95}$ (vence Z).
  \item \textbf{(B)} $\text{VP}_1 = \frac{44.000}{1{,}10} = 40.000$; $\text{VP}_2 = \frac{60.500}{1{,}21} = 50.000$; $\text{VP}_3 = \frac{66.550}{1{,}331} = 50.000$. Logo, $\text{VPL} = -100.000 + 40.000 + 50.000 + 50.000 = \mathbf{\reais{40.000}}$. Em $t=2$, o saldo acumulado é $-100.000 + 40.000 + 50.000 = -10.000$; portanto, $PB_d = 2 + \frac{10.000}{50.000} = \mathbf{2{,}20\text{ anos}}$.
  \item \textbf{(B)} $\text{WACC} = 0{,}70(15\%) + 0{,}30(10\%)(1 - 0{,}34) = 10{,}50\% + 1{,}98\% = \mathbf{12{,}48\%}$. $\text{Equity Value} = 1.500.000 - 450.000 + 120.000 = \mathbf{\reais{1.170.000}}$.
  \item \textbf{(B)} O desenvolvimento territorial exige convergência entre sustentabilidade financeira privada e retenção local de valor ($\text{IRV}$), encadeamento produtivo e governança ecossistêmica.
  \item \textbf{(B)} $N_{\text{AV3}} = 0{,}30(9{,}0) + 0{,}30(8{,}0) + 0{,}25(8{,}4) + 0{,}15(8{,}0) = 2{,}70 + 2{,}40 + 2{,}10 + 1{,}20 = \mathbf{8{,}40}$. $M_F = 0{,}30(7{,}5) + 0{,}20(8{,}0) + 0{,}50(8{,}40) = 2{,}25 + 1{,}60 + 4{,}20 = \mathbf{8{,}05}$.
  \item \textbf{(B)} A defesa oficial da AV3 segue o formato de 10 slides em 20 minutos mais 10 minutos de arguição, integrando todas as unidades da disciplina.
\end{enumerate}
\end{boxGabarito}

\begin{boxGabarito}[Espelho de Correção: Parte II (4 Questões Dissertativas)]
\begin{itemize}[leftmargin=1.4em, itemsep=4pt]
  \item \textbf{D1 (Unidade I):} Espera-se que o(a) discente explique que infraestrutura física isolada não gera inovação sem eficiência coletiva (economias externas incidentais e ações conjuntas deliberadas em APLs), sem distância cognitiva adequada (Boschma: nem excesso de proximidade que gere \textit{lock-in}, nem distância extrema que impeça aprendizado) e sem governança na Hélice Quádrupla/Quíntupla (incluindo sociedade civil e sustentabilidade).
  \item \textbf{D2 (Unidade II):} Espera-se articulação clara entre: (i) \textit{effectuation} na experimentação inicial com recursos do grupo de pesquisa e perda aceitável; (ii) quantificação realista de $\text{SOM}$ para validação comercial; (iii) registro de PI no INPI antes da exposição comercial e uso das regras de transferência de tecnologia das ICTs (Marco Legal de CT\&I, Lei 13.243/2016); e (iv) financiamento em estágios (\textit{grants} FAPESB/FINEP no Vale da Morte antes de crédito BNB/BNDES ou VC).
  \item \textbf{D3 (Unidade III):} Espera-se demonstração de que: (i) a $\text{TIR}$ pode divergir do $\text{VPL}$ em fluxos não convencionais ou projetos de escalas diferentes, além de não medir o $\text{IRV}$; (ii) a Matriz Locacional pondera externalidades e custos espaciais; (iii) o \textit{Payback} Descontado avalia o risco de liquidez no tempo; e (iv) o $\text{EV}$ representa o valor operacional total da firma, devendo-se descontar a dívida líquida para apurar o \textit{Equity Value} dos sócios.
  \item \textbf{D4 (Unidade IV):} Espera-se distribuição coerente nos 10 slides (Slides 1–2: $L_T$; Slides 3–6: $E_M$; Slides 7–9: $V_F$; Slide 10 e postura: $C_D$) e explicitação dos requisitos acadêmicos para publicação (pergunta de pesquisa, revisão de literatura, método auditável, fontes documentadas e análise crítica de limitações).
\end{itemize}
\end{boxGabarito}

\begin{tcolorbox}[colback=bgTeal, colframe=etTeal, arc=3pt, boxrule=0.8pt, title={\bfseries Painel de Diagnóstico Final, Prontidão para a AV3 e Publicação}, coltitle=white, colbacktitle=etTeal]
\renewcommand{\arraystretch}{1.2}
\footnotesize
\begin{tabularx}{\textwidth}{@{}c >{\raggedright\arraybackslash}X >{\raggedright\arraybackslash}X@{}}
  \toprule
  \rowcolor{etIndigo}
  \textcolor{white}{\bfseries Acertos (Parte I)} & \textcolor{white}{\bfseries Nível de Domínio} & \textcolor{white}{\bfseries Recomendação Pedagógica Imediata} \\
  \midrule
  \textbf{10 a 12} & \textbf{Domínio Pleno} & Apto(a) à defesa AV3 e à conversão do projeto em manuscrito para ENANPUR, ENANPAD ou ENGEMA. \\
  \rowcolor{bgGray}
  \textbf{7 a 9} & \textbf{Domínio Intermediário} & Revisar especificamente os capítulos correspondentes aos erros (ex.: $\text{HHI}$, $\text{LTV/CAC}$, $\text{VPL}$ ou $\text{WACC}$) no Formulário Técnico. \\
  \textbf{0 a 6} & \textbf{Reforço Necessário} & Refazer os Estudos de Caso (Etapa 4) e os Exercícios Propostos (Etapa 6) das unidades com maior déficit antes da banca AV3. \\
  \bottomrule
\end{tabularx}
\end{tcolorbox}
